\documentclass[10pt,a4paper]{amsart}
\usepackage[margin=19mm]{geometry}
\usepackage{amsmath,amssymb,mathtools}
\usepackage[scr=rsfs,cal=euler]{mathalfa}
\usepackage{xfrac}
\usepackage[unicode,hidelinks,bookmarks=false]{hyperref}
\newcommand{\ie}{i.e.\ }
\newcommand{\cf}{cf.\ }
\newcommand{\resp}{respectively\ }
\newcommand{\Subsection}{\subsection}
\providecommand{\nobreakdash}{\nobreak\hbox{-}}
\setlength{\emergencystretch}{2em}
\multlinegap0pt
\usepackage{bm}
\usepackage{bbm}
\usepackage{dsfont}
\usepackage{xspace}
\usepackage{cancel}
\usepackage{mathtools}
\usepackage{amsthm}
\makeatletter
\@ifundefined{theo}{\newtheorem{theo}{Theo}}{}
\@ifundefined{lem}{\newtheorem{lem}[theo]{Lemma}}{}
\@ifundefined{prop}{\newtheorem{prop}[theo]{Proposition}}{}
\@ifundefined{rem}{\newtheorem{rem}[theo]{Remark}}{}
\makeatother
\newcommand{\F}{\mathbb{F}}
\newcommand{\G}{\mathbb{G}}
\newcommand{\Ps}{\mathbb{P}}
\newcommand{\Spec}{\operatorname{Spec}}
\title[Hypergeometric Motives from Toric Hypersurfaces]{Hypergeometric Motives from Toric Hypersurfaces}
\author{Asem Abdelraouf}
\date{Source: arXiv:2509.17624v1 (CC BY 4.0)}
\begin{document}
\maketitle

\noindent\emph{Source and licence.} Hypergeometric Motives from Toric Hypersurfaces. Version arXiv:2509.17624v1 (CC BY 4.0), distributed under the Creative Commons Attribution 4.0 International licence, as recorded in the arXiv licence metadata for this exact version and shown on the abstract page https://arxiv.org/abs/2509.17624v1. Full rights notice: licenses/arxiv-2509.17624-CC-BY-4.0.txt.\par\smallskip

\noindent\textit{Editorial scope.} Counted unit: the Lem whose source label is \texttt{lem:fourier-series} in the article (source0.tex lines 1004--1014, source label \texttt{lem:fourier-series}), delivered together with the article's own complete original proof of it (source0.tex lines 1016--1082, one \texttt{proof} environment of 67 lines). No mathematics is written, repaired or re-derived: statement and proof are the article's own text line for line. Required context retained uncounted: lem:orthogonal (lines 398--411); lem:projective-to-toric (lines 832--842); rem:primitive-dehomogenization (lines 847--853); prop:main-count (lines 996--1002). Every definition, display and label of those lines is retained unchanged. Omitted from this package and not redistributed: 1--397, 412--831, 843--846, 854--995, 1003--1003, 1015--1015, 1083--2221. Nothing inside a retained excerpt is omitted or altered: every retained body is delivered exactly as the article's lines are written, so no omission receipt of any kind accompanies this package. Citations: the retained excerpts carry no citation command at all, so no bibliography entry of the article is transcribed and no reference is redistributed. Reference adaptation: none was needed and none was made; every reference inside the delivered unit resolves inside it, so no unresolved reference or citation is produced. Printed slips kept literal: no printed word, symbol, hypothesis or number of the counted statement or of its proof was altered. Layout and class: the article's own class is \texttt{amsart} with packages including graphicx, geometry, amsfonts, mathtools, amssymb, amsthm, booktabs, bm, mathabx, xcolor, float, longtable, hyperref, url; the wrapper is the lane's pinned \texttt{amsart} editorial wrapper at 10pt on A4, which restates the theorem-like environments the retained text needs and adds the article's own macro definitions that the retained text needs (\texttt{\textbackslash F}, \texttt{\textbackslash G}, \texttt{\textbackslash Ps}, \texttt{\textbackslash Spec}), with extra wrapper packages (bm, bbm, dsfont, xspace, cancel, mathtools). The delivered numbering is the unit's own. Assets: no retained span contains a figure environment, a graphics-inclusion command, a PDF-page inclusion, a table or a TikZ picture; no image file is copied into this package, the package declares no asset dependency and no image file is redistributed with it. Licence: version arXiv:2509.17624v1 of this article is distributed under the Creative Commons Attribution 4.0 International licence, as recorded in the arXiv licence metadata for this exact version and shown on the abstract page https://arxiv.org/abs/2509.17624v1. A full rights notice is delivered at licenses/arxiv-2509.17624-CC-BY-4.0.txt. Characters: every retained excerpt is pure ASCII, so the delivered engine, XeLaTeX with its default Latin Modern text font, reports no missing character.
\begin{lem} \label{lem:orthogonal} We have the following orthogonality relations:
\begin{equation*}
    \sum_{t \in \F_q} \psi(tx) = \begin{cases}
        q \quad & \text{ if } x =0,\\
        0 \quad & \text{ if } x\neq 0.
    \end{cases}
\end{equation*}
\begin{equation*}
   \sum_{m=0}^{q^\times-1} \chi^{m}(x) =   \begin{cases}
        q^\times \quad & \text{ if }x =1,\\
        0 \quad & \text{ if } x \neq 1.
        \end{cases}
\end{equation*}
\end{lem}

\begin{lem}\label{lem:projective-to-toric}
   If $Z$ is primitive, then it is isomorphic to the subvariety of $\Ps^{d+1}$ given by 
   \begin{align*}
    w_1+\dots +w_{d+2} &=0, \quad     w_1\cdots w_{d+2} \neq 0,\\
    \prod_{j=1}^{d+2} w_j^{\gamma_j} &=t,
    \end{align*}
    where  $w_1,\dots, w_{d+2}$ are homogeneous coordinates on $\Ps^{d+1}$ and $t = \prod_{j=1}^{d+2} u_j^{\gamma_j}$. Furthermore, such an isomorphism is given by 
    \[
    w_j = u_j \prod_{i=1}^{d} x_i^{m_{ij}},  \quad j=1,\dots, {d+2} .
    \]
\end{lem}

\begin{rem}\label{rem:primitive-dehomogenization}
The proof of Lemma \ref{lem:projective-to-toric} implies that the variety in $\Ps^{d+1}$ defined by $w_1\cdots w_{d+2} \neq 0$ 
and $\prod_{j=1}^{d+2} w_j^{\gamma_j} =t$ is isomorphic to $\G_m^{d} = \Spec k[x_1^{\pm},\dots,x_d^{\pm}]$ via the coordinate map
\[
w_j = u_j \prod_{i=1}^{d} x_i^{m_{ij}}, \quad j=1,\dots, d+2.
\]
\end{rem}

\begin{prop}\label{prop:main-count}
The number of $\F_q$-points of $W_S$ is given by
\begin{align*}
     \#W_{S} (\mathbb{F}_q) &= q^{-1}(q-1)^{d} + (-1)^{|S|}q^{-1}(q-1)^{|S|-1} \\
     &\cdot \sum_{\lambda \in \Lambda(q)}   \sum_{m=0}^{q^\times-1} \prod_{j\in S} \epsilon\left( -m \gamma_j -\delta(j,\lambda) \right) g\left(-m\gamma - \delta(\lambda)\right) \chi^{m}(t) \prod_{k=1}^d \chi^{\lambda_k}(\sigma_k).
\end{align*}
\end{prop}

\begin{lem}[Finite Fourier Transform]\label{lem:fourier-series}
Let $g: \left(\F_{q}^{\times}\right)^{d} \to \mathbb{C}$ be a complex valued function. Then, for all $x \in \left(\F_{q}^{\times}\right)^{d}$, we have
\[
g(x) = \sum_{i=1}^d \sum_{k_i =0}^{q^\times -1} a_k \prod_{i=1}^d \chi^{k_i}(x_i),
\]
where 
\[
a_k = \frac{1}{(q-1)^{d}}\sum_{x \in (\F_q^{\times})^{d} } g(x) \prod_{i=1}^d\chi^{-k_i}(x_i).
\]
Furthermore, the \emph{Fourier coefficients} $a_k$ are unique. 
\end{lem}

\begin{proof}[Proof of Proposition \ref{prop:main-count}]
Let $V$ be the subvariety in $\G_m^{d+2} \times \G_m^d$ defined by the equations
\begin{align}
f_0(w,z)&\coloneqq w^{\gamma} =t\eqqcolon \sigma_0, \label{eq:proof-main-count-1}\\
f_k(w,z)&\coloneqq w^{\rho_k} z^{-\mathcal{N}_k}=  \sigma_k, \qquad k=1,\dots,d.\label{eq:proof-main-count-2}
\end{align}
Note that $V$ is a cone over the variety $V'$ defined by the same homogeneous (in $w_1, \dots, w_{d+2}$) equations in $\Ps^{d+1} \times \G_m^d$. Setting 
\[
w_j = u_j \prod_{i=1}^{d} x_i^{f_{ij}}, \qquad j=1,\dots,d+2,
\]
Equation \eqref{eq:proof-main-count-1} becomes a tautology, and Equations \eqref{eq:proof-main-count-2} become:
\[
x_k z^{-\mathcal{N}_k}=  1, \qquad k=1,\dots,d.
\]
Thus, the coordinates $x_1,\dots,x_d$ are determined by $z_1,\dots,z_d$. Furthermore, the morphism defined by this substitution is an isomorphism by Remark \ref{rem:primitive-dehomogenization}. 
Thus, $V'$ is a torus of dimension $d$ and $V$ is a torus of dimension $d+1$.

We dehomogenise the equations of $W_S$ by setting $y_j = w_{j}/w_{d+2}$ for $j=1,\dots,d+1$. For $k=0,\dots,d$, we define $g_k$ by the equation
\[
g_k(y_1,\dots,y_{d+1},z)= f_k(y_1,\dots,y_{d+1},1,z).
\]
By Lemma \ref{lem:orthogonal}, we have:
\begin{align*}
     q  \# W_{S} (\mathbb{F}_q)  &= \sum_{
     \substack{w_{d+2}\in \F_q \\ y_1,\dots,y_{d+1},z_1,\dots,z_d \in \F_q^{\times} \\ g_k(y,z)=\sigma_k, \ k=0,\dots,d }} \psi\left(w_{d+2}\sum_{j \notin S} y_j\right) \\
     &= \sum_{
     \substack{y_1,\dots,y_{d+1},z_1,\dots,z_d \in \F_q^{\times} \\ g_k(y,z)=\sigma_k, \ k=0,\dots,d}} 1 + \sum_{
     \substack{w_{d+2}\in \F_q^\times \\ y_1,\dots,y_{d+1},z_1,\dots,z_d \in \F_q^\times \\g_k(y,z)=\sigma_k, \ k=0,\dots,d}} \psi\left(w_{d+2}\sum_{j \notin S} y_j\right).
\end{align*}
We make the change of variables $w_j = w_{d+2}y_j$,  $j=1,\dots, d+1$, then
\begin{align*}
      q  \# W_{S} (\mathbb{F}_q)  &= \sum_{
     \substack{y_1,\dots,y_{d+1},z_1,\dots,z_d \in \F_q^{\times} \\ g_k(y,z)=\sigma_k, \ k=0,\dots,d}} 1 + \sum_{
     \substack{w_1,\dots, w_{d+2}, z_1, \dots, z_d \in \F_q^\times \\f_k(w,z)=\sigma_k, \ k=0,\dots,d}} \psi\left(\sum_{j\notin S} w_j\right)\!.
\end{align*}
The first term is the number of points of the dehomogenisation of the equations defining $V$, thus it is the number of points of a torus of dimension $d$. Thus, 
\begin{align*}
       q  \# W_{S} (\mathbb{F}_q)  &= (q-1)^{d} + \sum_{
     \substack{w_1,\dots, w_{d+2}, z_1, \dots, z_d \in \F_q^\times \\f_k(w,z)=\sigma_k, \ k=0,\dots,d}} \psi\left(\sum_{j\notin S} w_j\right)\!\\
    &= (q-1)^{d} + \sum_{(w,z) \in V(\F_q) } \psi\left(\sum_{j \notin S} w_j\right).
    \end{align*}
The last sum depends on $t, \sigma_1,\dots, \sigma_d$, so we consider $q \# W_{S}(\F_q)$ as a function of variables $t, \sigma_1,\dots, \sigma_d$ and compute its Fourier expansion.
By Lemma \ref{lem:fourier-series}, we have 
\begin{align*}
 q \# W_{S}(\F_q) = (q-1)^{d} + \frac{1}{ (q-1)^{d+1}}\sum_{(m,\lambda) \in \left(\F_q^{\times}\right)^{d+1}} c_{m,\lambda} \cdot   \chi^{m}(t) \prod_{k=1}^d \chi^{\lambda_k} (\sigma_k),
\end{align*}
   where 
\begingroup
\allowdisplaybreaks
\begin{align*}
    c_{m,\lambda}  &= \sum_{(t,\sigma) \in \left(\F_q^\times\right)^{d+1}} \sum_{(w,z) \in V(\F_q)} \psi \left(\sum_{j\notin S} w_j\right) \chi^{-m}(t) \prod_{k=1}^d \chi^{-\lambda_k} (\sigma_k) \\
    &=  \sum_{(w,z) \in \left(\F_q^\times\right)^{2d+2} } \psi \left(\sum_{j\notin S} w_j\right) \chi^{-m}(w^\gamma) \prod_{k=1}^d \chi^{-\lambda_k}\left( w^{\delta_k}z^{-\mathcal{N}_k}\right) \\
    &=  \prod_{j \notin S}  \sum_{w_j \in \F_q^\times} \psi(w_j) \chi^{-m \gamma_j -\sum_{k=1}^d \lambda_k\delta_{kj}}(w_j) \cdot \prod_{j \in S}  \sum_{w_j \in \F_q^\times} \chi^{-m \gamma_j -\sum_{k=1}^d \lambda_k\delta_{kj}}(w_j) \\
    &\cdot  \prod_{i=1}^d \sum_{z_i \in \F_q^\times} \chi^{ - \sum_{k=1}^d \lambda_k\mathcal{N}_{ki}}(z_i) \\
    &=  \prod_{j \notin S}   g\left(-m\gamma_j -  \sum_{k=1}^d \lambda_k \delta_{kj}\right) \cdot \prod_{j\in S} (q-1)\epsilon\left( -m \gamma_j -\sum_{k=1}^d \lambda_k\delta_{kj} \right) \\
    &\cdot  \prod_{i=1}^d (q-1) \epsilon \left( \sum_{k=1}^d \lambda_k\mathcal{N}_{ki} \right) \\
    &= (-1)^{|S|} (q-1)^{d+|S|}\cdot \prod_{j=1}^{d+2} g\left(-m\gamma_j -  \sum_{k=1}^d \lambda_k \delta_{kj}\right)   \\
    &\cdot\prod_{j\in S} \epsilon\left( -m \gamma_j -\sum_{k=1}^d \lambda_k\delta_{kj} \right) \cdot \prod_{i=1}^d \epsilon \left( \sum_{k=1}^d \lambda_k\mathcal{N}_{ki} \right).
\end{align*}
\endgroup
The product $\prod_{i=1}^d \epsilon \left( \sum_{k=1}^d \lambda_k\mathcal{N}_{ki} \right)$ is $1$ if $\lambda \in \Lambda(q)$, and is $0$ if $\lambda \notin \Lambda(q)$.  Thus, we have 
\begin{align*}
    q \# W_{S} (\mathbb{F}_q) &= (q-1)^{d} + (-1)^{|S|}(q-1)^{|S|-1}\\
    &\cdot\sum_{\lambda \in \Lambda(q)}  \sum_{m=0}^{q^\times -1}\prod_{j\in S} \epsilon\left( -m \gamma_j -\delta(j,\lambda) \right)\prod_{j=1}^{d+2}g\left(-m\gamma_j - \delta(j,\lambda)\right) \chi^{m}(t) \chi^{\lambda}(\sigma),
\end{align*}
as desired.
\end{proof}

\end{document}
